\documentclass[10pt,a4paper]{amsart}
\usepackage[margin=19mm]{geometry}
\usepackage{amsmath,amssymb,mathtools}
\usepackage[scr=rsfs,cal=euler]{mathalfa}
\usepackage{xfrac}
\usepackage[unicode,hidelinks,bookmarks=false]{hyperref}
\newcommand{\ie}{i.e.\ }
\newcommand{\cf}{cf.\ }
\newcommand{\resp}{respectively\ }
\newcommand{\Subsection}{\subsection}
\providecommand{\nobreakdash}{\nobreak\hbox{-}}
\setlength{\emergencystretch}{2em}
\multlinegap0pt
\usepackage{amssymb}
\usepackage{mathtools}
\usepackage{dsfont}
\usepackage[shortlabels]{enumitem}
\usepackage{float}
\usepackage{xcolor}
\newtheorem{lemma}{Lemma}
\newtheorem{theorem}{Theorem}
\title{A note on idempotents in quandle rings}
\author{Dilpreet Kaur, Pushpendra Singh}
\date{Source: arXiv:2609.09183v1 (CC BY 4.0)}
\begin{document}
\maketitle

\noindent\emph{Source and licence.} A note on idempotents in quandle rings. Version arXiv:2609.09183v1 (CC BY 4.0), distributed by arXiv under the Creative Commons Attribution 4.0 International licence, http://creativecommons.org/licenses/by/4.0/, as recorded in the arXiv licence metadata for this exact version and shown on the abstract page https://arxiv.org/abs/2609.09183v1. Full licence notice: \texttt{licenses/arxiv-2609.09183-CC-BY-4.0.txt}.\par\smallskip

\noindent\textit{Editorial scope.} Counted unit: the article's theorem theorem (source0.tex lines 286--288). The article's complete original proof of it is delivered with it (source0.tex lines 290--316, one \texttt{proof} environment of 27 lines). No mathematics is written, repaired or re-derived: the statement and the proof are the article's own lines, in the article's own notation, and no retained line is substituted or re-flowed.  Retained uncounted as the source-local context the proof needs: lines 161--166.  Nothing was omitted from any retained span, so the package carries no omission record.  No retained line contains a citation, so the wrapper carries no bibliography and no bibliography entry.  No retained line contains a figure environment, an image-inclusion command or drawing code, so no image is delivered and the package declares no asset dependency.  Editorial formatting only, in addition: an AMS \texttt{amsart} preamble that restates the article's own declarations and macros used by the retained lines, namely the environments \texttt{lemma}, \texttt{theorem} on separate counters, together with the standard packages those lines need.  The retained environments print in this document as Lemma 1, Theorem 1 in order of appearance, whereas the article prints the same items with the numbering of its own full document.  The complete native source of the article as distributed in the arXiv e-print for this exact version is bundled unchanged as \texttt{sources/209787/source0.tex}.
\begin{lemma}\label{inv}
Let $x\in \mathcal R_n$. Then we have
\[a_x= \frac{1}{n} \sum\limits_{r\in \mathbb Z_n} \hat{a}(r)\zeta^{-rx}\]
Furthermore,
\[ \sum\limits_{r\in \mathbb Z_n} |\hat{a}(r)|^2=n \sum\limits_{x\in \mathbb Z_n} a_x^2 \]
\end{lemma}

\begin{theorem}
Let the quandle ring $\mathbb Z[Q]$ for a connected quandle $Q={\rm Aff}(\mathbb Z_p,f)$ have a nontrivial idempotent $u$ with augmentation $1$. Then $\mathbb Z[Q]$ has infinitely many nontrivial idempotents.
\end{theorem}

\begin{proof}
Let $f(x)=tx$ for $t\in \mathbb Z_{p}^{*}$, then the quandle operation is $x\triangleright y=tx+(1-t)y$. Let $s=1-t$, then $s\neq 0$, since $Q$ is connected.  Let $u = \sum\limits_{x\in \mathbb Z_p}a_xe_x\in \mathbb Z[Q]$ be a nontrivial idempotent with augmentation $1$. Let $\zeta=e^{i2\pi/p}$ and define $\hat{a}(r)=\sum\limits_{x\in \mathbb Z_p}a_x\zeta^{rx}$ for all $r\in \mathbb Z_{p}$. We have $u^2=\sum\limits_{x,y\in \mathbb Z_p}a_xa_ye_{tx+sy}$ and so the coefficient of $e_z$ in $u^2$ is $\sum\limits_{tx+sy=z}a_xa_y$.



Construct the associative ring $R=\mathbb Z[X]/\langle X^{p}-1 \rangle$. Using the coefficients of the nontrivial idempotent $u$, define an element of $R$ as $A(X)=\sum\limits_{x=0}^{p-1}a_xX^x$. We have 
\[ A(X^{t})A(X^{s})= \sum\limits_{x,y\in \mathbb Z_p}a_xa_yX^{tx+sy}\]
The coefficients of $X^z$ modulo $X^{p}-1$ in the above expression
are exactly the coefficients of $e_z$ in $u^2$. Thus $u^2=u$ in $\mathbb Z[Q]$ is equivalent to $A(X^{t})A(X^{s})=A(X)$ in $R$.

Now for all positive integers $m\geq 1$, define $A_m(X)=A(X)^m$ and write $A_m(X)=\sum\limits_{z=0}^{p-1}a_z^{(m)}X^z \in R$. Using these coefficients, define $u_m=\sum\limits_{z \in \mathbb Z_p}a_z^{(m)}e_z$. Now we show that $u_m$ is an idempotent. It is enough to show $A_m(X^{t})A_{m}(X^{s})=A_m(X)$. We have
$A_m(X^{t})A_m(X^{s})=A(X^{t})^mA(X^{s})^m=(A(X^{t})A(X^{s}))^m=A(X)^m=A_m(X)$. Hence $u_m$ is an idempotent.

Now we show that $u_m\neq 0$. We have $\varepsilon(u)=\sum\limits_{x\in \mathbb Z_p}a_x=A(1)=1$. Now $\varepsilon(u_m)=\sum\limits_{x\in \mathbb Z_p}a_x^{(m)}=A_m(1)=A(1)^m=1^m=1$.

Now we show that for all $m\geq 1$, $u_m$ are distinct. Define a ring homomorphism $\phi: \mathbb Z[X]/\langle X^{p}-1 \rangle \rightarrow \mathbb C, F(X)\mapsto F(\zeta^r)$. We note that we can always choose $0\neq r\in \mathbb Z_{p}$ such that $|A(\zeta^r)|\neq 1$. Because if not, then $|A(\zeta^r)|=|\hat{a}(r)|=1$ for all $r\neq 0$, and since $A(\zeta^0)=A(1)=1$. Then Lemma \ref{inv} implies $\sum\limits_{x\in \mathbb Z_p}a_x^2=1$ and so $u$ is trivial, which is a contradiction. 
Moreover, we have $A(\zeta^r)\neq 0$, because if $A(\zeta^r)=0$, then $A(X)=c(1+X+\dots+X^{p-1})$ for some $c\in \mathbb Z$ and so $A(1)=cp\neq 1$.


Let $\rho=|A(\zeta^r)|$. Then $\rho>0$ and $\rho\neq 1$. Now $|A_m(\zeta^r)|=|A(\zeta^r)|^m=\rho^m$. Now for $m\neq k$, we get $\rho^m\neq \rho^k$ and so $A_m(\zeta^r)\neq A_k(\zeta^r)$ and consequently $A_m(X)\neq A_k(X)$. Thus $u_m\neq u_k$.


Finally, suppose $u_m=e_j$ for some $j\in \mathbb Z_p$. Then $A_m(X)=X^j$ and so $|A_m(\zeta^r)|=|\zeta^{rj}|=1$ but we have $|A_m(\zeta^r)|=|A(\zeta^r)|^m=\rho^m\neq 1$. Thus $u_m,m\geq 1$ are nontrivial, augmentation $1$, distinct idempotent elements of $\mathbb Z[Q]$.



\end{proof}

\end{document}
